\chapter{Complexity — Meyer Sigma}
%%% generated by the blueprint pipeline from TCSlib.Complexity.CircuitComplexity.MeyerSigma
\section{Overview}

Declarations of the Lean module \texttt{TCSlib.Complexity.CircuitComplexity.MeyerSigma}.

\section{Declarations}

\begin{definition}[Bounded binary word enumeration]
\lean{Complexity.Meyer.wordOf}
\label{Complexity.Meyer.wordOf}
\leanok
\uses{Complexity.enumWord}
For natural numbers $W$ and $v$, let $\text{enumWord}(W,v)$ denote the binary word
enumerating $v$ at width $W$. Define the associated word to be $\text{enumWord}(W,v)$ if
$v<2^W$, and the empty word otherwise.
\end{definition}

\begin{definition}[Answer from a Bounded Table]
\lean{Complexity.Meyer.AnsOf}
\label{Complexity.Meyer.AnsOf}
\leanok
\uses{Complexity.Meyer.Sel, Complexity.Meyer.ansD, Complexity.Meyer.wordOf}
Let $M$ be a finite Turing machine over the Boolean alphabet, let $W$ be a natural
number, let $x,D$ be Boolean strings, and let $\text{sel}$ be a selector. For natural
numbers $t,a$, define the answer by applying the table-answer procedure to
$M,D,\text{sel},x$, and the length-$W$ enumerations of $t$ and $a$, using the empty word
whenever either number is at least $2^W$.
\end{definition}

\begin{definition}[Membership Bit for a Circuit Value]
\lean{Complexity.Meyer.xbOf}
\label{Complexity.Meyer.xbOf}
\leanok
\uses{BoolCircuit.CVAL, Complexity.Meyer.idxDesc, Turing.pairEncode}
Let $M$ be a finite Turing machine over the Boolean alphabet, let $x$ be a Boolean
string, and let $j$ be a natural number. The associated bit is $1$ exactly when the pair
encoding of the description of index $j$ for strings of length $|x|$, together with $x$,
belongs to the circuit-value language; otherwise it is $0$.
\end{definition}

\begin{definition}[Circuit for an Input Index]
\lean{Complexity.Meyer.idxCircuit}
\label{Complexity.Meyer.idxCircuit}
\leanok
\uses{BoolCircuit.DAGCircuit}
For natural numbers $n,j$ with $j<n$, the circuit with $n$ inputs and no gates whose
output is input vertex $j$ is well-defined.
\end{definition}

\begin{theorem}[Encoding an input-index circuit]
\lean{Complexity.Meyer.idxCircuit_encode}
\label{Complexity.Meyer.idxCircuit_encode}
\leanok
\uses{BoolCircuit.DAGCircuit.encode, BoolCircuit.encodeList, BoolCircuit.encodeNat, Complexity.Meyer.idxCircuit, Complexity.Meyer.idxDesc}
\difficulty{1}
For natural numbers $n,j$ with $j<n$, the Boolean-word encoding of the circuit on $n$
inputs with no gates and output vertex $j$ is the word consisting of $n$ true bits
followed by false, then two false bits, then $j$ true bits followed by false.
\end{theorem}

\begin{theorem}[Recovered bit from a circuit encoding]
\lean{Complexity.Meyer.xbOf_eq}
\label{Complexity.Meyer.xbOf_eq}
\leanok
\uses{BoolCircuit.CVAL, BoolCircuit.DAGCircuit.eq_of_encode_eq, BoolCircuit.DAGCircuit.eval, BoolCircuit.DAGCircuit.values, Complexity.Meyer.idxCircuit, Complexity.Meyer.idxCircuit_encode, Complexity.Meyer.idxDesc, Complexity.Meyer.xbOf, Complexity.pairFstD, Complexity.pairFstD_pairEncode, Complexity.pairSndD, Complexity.pairSndD_pairEncode, Turing.pairEncode}
\difficulty{5}
Let $x$ be a Boolean word and let $j$ be a natural number with $j<|x|$. The bit
indicating whether the pair encoding of the index description for $(|x|,j)$ followed by
$x$ belongs to the language of well-formed fan-in-two Boolean circuits that evaluate to
true is exactly the $j$th bit of $x$.
\end{theorem}

\begin{theorem}[Bounded Word Enumeration]
\lean{Complexity.Meyer.wordOf_of_lt}
\label{Complexity.Meyer.wordOf_of_lt}
\leanok
\uses{Complexity.Meyer.wordOf, Complexity.enumWord}
\difficulty{1}
For natural numbers $W$ and $v$ with $v<2^W$, the word assigned to $v$ at width $W$ is
its width-$W$ little-endian binary word, with high-order zeros retained.
\end{theorem}

\begin{theorem}[Atom valuations on a valid encoded instance]
\lean{Complexity.Meyer.atoms_vin}
\label{Complexity.Meyer.atoms_vin}
\leanok
\uses{Complexity.KarpLipton.dropMarker, Complexity.Meyer.AnsOf, Complexity.Meyer.NumVal, Complexity.Meyer.avOf, Complexity.Meyer.awd, Complexity.Meyer.bits_pad, Complexity.Meyer.bjw, Complexity.Meyer.ctrVal_enumWord, Complexity.Meyer.ctrVal_lt, Complexity.Meyer.enumWord_ctrVal, Complexity.Meyer.enumWord_ones, Complexity.Meyer.enumWord_zero, Complexity.Meyer.evalAtom, Complexity.Meyer.incFixed_enumWord, Complexity.Meyer.incFixed_eq_none_iff, Complexity.Meyer.tvOf, Complexity.Meyer.twd, Complexity.Meyer.vD_vin, Complexity.Meyer.va_vin, Complexity.Meyer.vin, Complexity.Meyer.vj_vin, Complexity.Meyer.vp_vin, Complexity.Meyer.vt_vin, Complexity.Meyer.vx_vin, Complexity.Meyer.wd, Complexity.Meyer.wordOf, Complexity.Meyer.wordOf_of_lt, Complexity.Meyer.xbOf, Complexity.TimeHierarchy.ctrVal, Complexity.enumWord, Turing.incFixed}
\difficulty{6}
Let $C_w,c_w$ be natural numbers, and let $x,U,t_w,a_w,j_w,p_w$ be Boolean words.
Suppose $|x|+2<2^{C_w(|x|+1)^{c_w}}$ and the length-validity atom evaluates to true on
the nested pair encoding of these words. Then the function assigning each atom its
evaluation on that encoding is exactly the numerical valuation obtained from the encoded
time and auxiliary counters, the input length and index, the marked word derived from
$U$, and the input bit predicate, with answer atoms evaluated by the bounded-time answer
function.
\end{theorem}

\begin{theorem}[Circuit Query Membership Characterization]
\lean{Complexity.Meyer.query_mem_CVAL_iff}
\label{Complexity.Meyer.query_mem_CVAL_iff}
\leanok
\uses{BoolCircuit.CVAL, BoolCircuit.DAGCircuit.encode, BoolCircuit.DAGCircuit.encode_injective, BoolCircuit.DAGCircuitFamily, BoolCircuit.DAGCircuitFamily.HasFaninTwo, BoolCircuit.DAGCircuitFamily.language, BoolCircuit.DAGCircuitFamily.mem_language_iff, Turing.pairEncode, Turing.pairEncode_injective}
\difficulty{4}
Let $F$ be a family of Boolean circuits such that every circuit in the family is well
formed and has fan-in at most two, and let $q$ be a bit string. The encoding of the
circuit in the family for inputs of length $|q|$, paired with $q$, belongs to the
language of satisfiable circuit encodings if and only if that circuit accepts $q$.
\end{theorem}

\begin{theorem}[Circuit family computes bounded simulation queries]
\lean{Complexity.Meyer.ansD_family}
\label{Complexity.Meyer.ansD_family}
\leanok
\uses{BoolCircuit.DAGCircuit.encode, BoolCircuit.DAGCircuitFamily, BoolCircuit.DAGCircuitFamily.HasFaninTwo, BoolCircuit.DAGCircuitFamily.language, Complexity.Meyer.AnsOf, Complexity.Meyer.Sel, Complexity.Meyer.Tab, Complexity.Meyer.ansD, Complexity.Meyer.ctrVal_enumWord, Complexity.Meyer.encodeSel, Complexity.Meyer.nbits, Complexity.Meyer.query, Complexity.Meyer.query_mem_CVAL_iff, Complexity.Meyer.tabAns, Complexity.Meyer.tabFun, Complexity.Meyer.tabFun_query, Complexity.Meyer.wordOf_of_lt, Complexity.enumWord, Turing.MultiTapeTM.initCfg, Turing.MultiTapeTM.runFrom}
\difficulty{4}
Let $M$ be a finite-state Turing machine with Boolean alphabet, and let $F$ be a family
of Boolean circuits such that every circuit has fan-in at most two and $F$ accepts
exactly the language recognized by the tabulating function of $M$. For any natural
number $W$, Boolean word $x$, selector $\text{sel}$, and natural numbers $t,a<2^W$, the
circuit in $F$ on $n_{
m bits}(M)+4W+4+|x|$ input bits, encoded as a circuit description
and queried with the encoding of $\text{sel}$ and the width-$W$ words for $t$ and $a$
together with $x$, returns the tabulated answer for $\text{sel}$ after $t$ steps of $M$
on $x$, with counter value $a$.
\end{theorem}
